% ITI 2026 (NIT Agartala, 13-14 Nov 2026) -- Springer LNNS, 10-12 pages.
%
% Build:  pdflatex iti2026 && bibtex iti2026 && pdflatex iti2026 && pdflatex iti2026
%
% MANDATED STRUCTURE (ITI 2026 rules): abstract, keywords, introduction,
% literature review, methodology, results, discussion, conclusion, references.
% All nine must be present. This is NOT the draft's structure -- the draft has
% threat model / system / decision rule / calibration / consistency / evaluation
% / ablations / limitations / related work as top-level sections. The remap:
%
%   introduction        <- 01-introduction
%   literature review   <- 11-related-work           (moves UP to section 2)
%   methodology         <- 02 + 03 + 04 + 05 + 06 + 07   (six sections merged)
%   results             <- 08-evaluation (+ 1 sentence of 09-ablations)
%   discussion          <- interpretation + 10-limitations
%   conclusion          <- 12-conclusion
%
% The consequence worth stating: the decision rule -- the contribution -- is no
% longer a top-level section. It becomes 3.2, and gets the largest subsection
% budget in the paper. Do not let the merge bury it.
%
% BUDGET (see writing/paper/CUT-PLAN.md). 12 pages leaves ~7.6 pages of body text
% => ~4,900 words. The markdown draft is 20,138. Each section carries its budget
% in a comment; if a section is over, something has to go, not the margins.
%
\documentclass[runningheads]{llncs}

\usepackage[T1]{fontenc}
\usepackage{graphicx}
\usepackage{booktabs}
\usepackage{amsmath}
\usepackage{amssymb}   % \mathbb for the expectation operator
\graphicspath{{../../figures/pdf/}}

\begin{document}

\title{Pricing Deception: When a Web-Attack Detector\\Should Probe Rather Than Decide}
\titlerunning{Pricing Deception}

\author{Parth Kansal \and Amrit Singh Nijjar \and Siddhant Mehta}
\authorrunning{P. Kansal et al.}
% ITI requires the affiliation to carry only department, institution, city and
% country -- no post or position.
\institute{Department of AIT-CSE, Chandigarh University,\\
Gharuan, Mohali, Punjab, India\\
\email{parth.kansal823@gmail.com, amrit.nijj814@gmail.com, 23bis70162@cuchd.in}}

\maketitle

% ================================================================ abstract 180w
\begin{abstract}
A web application firewall must commit, allow or block, on the evidence one request
carries. Deception normally happens after that commitment: a session already judged
hostile is moved into a honeypot. We treat it instead as a move available while the
detector is still unsure. When intent is uncertain the defence adds an invisible,
inert probe to the response, such as a fake table named in an error. An honest
client never renders it; someone probing the application acts on it, and reveals
intent. Our contribution is the rule deciding when to deploy one. A probe brings no
immediate benefit, since the request still reaches the real application, so its
whole worth is the information a bite reveals. Priced as the expected value of that
information, probing becomes cost-optimal over a belief band whose edges follow from
a frozen cost table and a measured bite likelihood ratio. Under cost accounting
alone the band is empty, so there is no middle action to tune. Over 99 seeded
traffic draws against one frozen model, a randomised holdout puts the probe's causal
effect at $+0.070$ $[+0.052, +0.088]$ (Fisher exact, $p = 3.4 \times 10^{-19}$), and
recall rises from 0.889 to 0.943 (paired McNemar, $p = 1.9 \times 10^{-95}$),
concentrated in the one attack category the passive detector cannot settle. No
legitimate session was diverted, including the four the passive baseline lost.

\keywords{cyber deception \and value of information \and cost-sensitive detection
\and web application security \and honeytokens.}
\end{abstract}

% ============================================================ 1 introduction 600w
\section{Introduction}

A web application firewall makes a binary choice on every request: let it through or
block it, on what that request alone reveals. The trade-off is familiar. Act early
on thin evidence and legitimate users get blocked; wait for evidence strong enough
to be safe and the attacker has had several requests' head start. Learned detectors
soften the edges but do not escape it. They still watch, wait, and commit on one
accumulated score.

Deception is the usual answer to passivity, yet in almost every deployment it comes
\emph{after} the decision: a request is judged hostile, and only then is the
session moved into a honeypot. Deception is a consequence of detection, not an
instrument of it. We have not found prior work treating it as a move available
while the detector is still unsure.

That move is the subject of this paper. When the system is uncertain, it adds
something to the response that a browser never renders but anyone reading raw
bytes will see: a fake database error, an unused JSON field, a hint at a
deprecated endpoint. An honest user notices nothing; someone probing the
application acts on it, and reveals what a passive score could only wait for. We
call these probes \emph{bait}. Our question is not how to build them, since
invisibility and consistency are engineering, but when to deploy one at all.

That decision need not be a heuristic. A probe brings no immediate benefit, since
the request it rides on still reaches the real application, so baiting an attacker
costs what letting one through costs. Its whole value is the information a bite
reveals, and that value can be computed. Priced as the expected value of the sample
information it buys, a probe becomes the cost-optimal action over a belief band
\emph{derived} from the cost of the defender's errors and the \emph{measured}
effectiveness of the probe. Account for the same costs without the information term
and the band vanishes: there is nothing there to tune.

\paragraph{Contributions.}
The ingredients are old. New here is the rule deciding \emph{when} to deceive, and
that its parameters are derived rather than chosen.

\begin{enumerate}
\item \textbf{A priced third action} (Section~\ref{sec:rule}). Probing maximises the
expected value of sample information over a band whose edges come out of a frozen
cost table and calibrated probe effectiveness. Under cost accounting alone that band
is empty, so the middle action is priced rather than tuned.

\item \textbf{Bait as a randomised treatment} (Section~\ref{sec:results}). A fixed
fraction of sessions reaching the band are deliberately left unbaited. Assignment is
random at the same belief state, so the gap between the arms estimates the
\emph{causal} effect of baiting instead of comparing two systems.

\item \textbf{Consistency across the divert boundary}
(Section~\ref{sec:consistency}). A decoy can be perfectly self-consistent while
contradicting everything the attacker saw before being moved. We show that
self-consistency is the wrong estimand, and measure the right one.
\end{enumerate}

We test against off-the-shelf attack tools we did not write and against the OWASP
ModSecurity Core Rule Set on identical traffic, and we report where the approach
does not win (Section~\ref{sec:discussion}).

% ======================================================= 2 literature review 450w
\section{Literature Review}
\label{sec:related}

Every ingredient here is established. What is new is where the decision to deceive
sits and how its parameters are set.

\paragraph{Detection is passive, and blind where syntax is valid.}
Signature firewalls match known-bad patterns and are blind to attacks that are valid
syntax. Learned detectors trade brittleness for a model, from early anomaly-based
detectors of web requests~\cite{kruegel2003anomaly} to recent adaptive
approaches~\cite{amouei2022rat}. All share the passive stance: watch a session
accumulate evidence, then commit on a score. Our detector is one of these and is not
the contribution; a signature firewall and the passive meter are our baselines
because the contribution sits on top of them. A separate line asks whether a client
is a script rather than whether it is hostile~\cite{iliou2019towards}, and its most
useful finding for us is negative: advanced bots imitate browser fingerprints closely
enough that automation signals alone stop separating them from people.

\paragraph{Deception is a destination, not an instrument.}
Deception as defence is well studied, from virtual honeypot
frameworks~\cite{provos2004honeyd} through game-theoretic
treatments~\cite{pawlick2019taxonomy} to whether decoys actually change attacker
behaviour~\cite{fergusonwalter2021examining}. In almost all of it deception is a
\emph{destination}: a caught attacker is moved somewhere, after a decision already
taken. Planting a credential as a tripwire is likewise old and widely
deployed~\cite{yuill2004honeyfiles,juels2013honeywords}, but that literature treats
honeytokens as always-on. Recent application-layer work automates their deployment,
rotation and teardown properly~\cite{kahlhofer2024applayer,kahlhofer2025koney} and
stops where we begin: it does not decide \emph{when} to deploy one from a belief
about the visitor being served.

\paragraph{Generated decoys are measured for fidelity, not contradiction.}
Language models increasingly generate convincing interactive
honeypots~\cite{sladic2024shellm,bridges2025sok}. That line is complementary rather
than competing, since our consistency layer is generator-agnostic and a language
model is exactly what it is built to sit in front of. What it lacks is the
guarantee: evaluation frameworks measure stealth and
fidelity~\cite{vero2026honeyval}, not whether a decoy contradicts itself over an
engagement --- and they measure it \emph{within} the decoy, which
Section~\ref{sec:consistency} argues is the wrong estimand for a decoy an attacker
is moved into.

\paragraph{Cost-sensitive detection has no information-purchase action.}
Our machinery is decades old: the expected value of sample
information~\cite{howard1966information}, cost-sensitive decisions under a loss
matrix~\cite{elkan2001foundations}, signalling models of deception with
evidence~\cite{pawlick2019leaky}, and the base-rate constraints binding any
intrusion detector~\cite{axelsson2000baserate}. None offers an action that buys
evidence, so none can express \emph{probe in order to find out} --- the gap
Section~\ref{sec:rule} fills. Throughout we follow the methodological cautions
of~\cite{sommer2010outside,arp2022dos}.

% ============================================================= 3 methodology 2100w
\section{Methodology}
% SOURCE 02 + 03 + 07 + 04 + 05 + 06 (9,940w -> 2,100w). The largest merge in the
% paper. Four subsections; 3.2 is the contribution and gets the most space.

\subsection{System and Threat Model}
\label{sec:system}

We consider an attacker whose intent cannot be settled from any single request. Some
requests carry obfuscated payloads that a signature written for the canonical form
misses~\cite{amouei2022rat}; others carry no payload at all. The clearest case, and
the one we lean on throughout, is an insecure direct object reference:
\texttt{/records/5} is syntactically indistinguishable from \texttt{/records/6}, yet
if the second object is not yours the request is an attack and nothing in its bytes
says so. There is no signature to write, and a behavioural detector struggles because
a benign integration reading its own records in order produces the same shape. This
is where a passive score is genuinely uncertain.

Attackers vary on two independent axes, automation and intent. A scanner is
automated and hostile, a price-comparison bot automated and harmless, a careful
human manual and hostile. Modern bots imitate browser fingerprints well enough that
automation signals alone stop separating them from people~\cite{iliou2019towards},
so collapsing both axes onto one score cannot tell the second case from the third.
Half of our simulated attack sessions therefore drive a real browser and fetch page
sub-resources, since much current tooling is built on Selenium or Playwright and
inherits that whatever the operator intends. Out of scope are network-layer attacks,
denial of service, and client-side attacks on other users. We assume the attacker
has not previously mapped the decoy world; honeypots have been fingerprinted at
scale~\cite{vetterl2018bitter}.

The defender runs a reverse proxy in front of the real application, which never sees
an unfiltered request nor knows the deception machinery exists. Every request
follows the same path: identify the session by cookie, with a fingerprint fallback
for clients that drop them; reduce it to eighteen features computed only from what a
live proxy can observe; update a two-part suspicion score; price the three actions
and take the cheapest; write the decision to an append-only, hash-chained log.

The features split along the two axes. Ten \emph{automation} features describe how
the client behaves: inter-request timing and its regularity, whether static assets
are fetched, browser-header completeness, whether a cookie is carried. Eight
\emph{malice} features describe what it is trying to do: special-character density
in client-controlled input, database keywords in syntactic context, error-response
ratio, failed authentications, and how many distinct accounts it has tried to log in
as. Two logistic heads, one per axis, accumulate across the session. Separating them
earns its keep twice: cost depends on hostility alone, so the belief driving the
policy is the malice score, while the automation score is spent on choosing
\emph{which} probe to deploy, because a scanner and a careful human take different
bait at different rates.

% architecture.pdf removed: the paragraph above walks the same pipeline in order,
% so the figure restated it and cost ~0.4 page against a hard 12-page limit.

\subsection{Pricing the Probe as an Information Purchase}
\label{sec:rule}

At each point in a session the defender holds a belief $p \in [0,1]$ that the
client is hostile, and chooses one of three actions. \emph{Pass} forwards the
request to the real application. \emph{Divert} moves the session into the decoy.
\emph{Bait} injects an invisible, inert probe into the response and forwards the
request unchanged. Each (action, true-class) pair carries a cost fixed in
Table~\ref{tab:costs}, which is hashed and frozen before any data is collected, so
no threshold derived below can have been tuned to the results.

Three entries carry the argument. Diverting a benign user must almost never happen,
so it is priced two orders of magnitude above any other error. Baiting a hostile
client costs exactly what passing one costs, because the request still reaches the
real application, so the only difference between passing and baiting is the small
nuisance cost borne by the benign probability mass. Taking expectations under $p$,
each action's immediate expected cost is affine:
%
\begin{equation}
\mathbb{E}[C_{\text{pass}}] = 25p, \quad
\mathbb{E}[C_{\text{bait}}] = 1 + 24p, \quad
\mathbb{E}[C_{\text{divert}}] = 200 - 220p .
\end{equation}

\paragraph{Under cost accounting alone there is no third action.}
Comparing the first two, $\mathbb{E}[C_{\text{bait}}] - \mathbb{E}[C_{\text{pass}}]
= 1 - p \geq 0$ for every belief: baiting is \emph{strictly} more expensive than
passing everywhere, by exactly the residual benign-nuisance term. The middle action
is therefore never chosen, and the policy collapses to a two-way rule with a single
crossover at $25p = 200 - 220p$, i.e. $p^{*} = 200/245 \approx 0.816$ --- one
PASS/DIVERT boundary and \textbf{no band anywhere in between}. This is the sharpest
answer to the natural objection that the middle action is a tuned threshold:
\emph{there is no middle action to tune}. It appears only once the information a
probe buys is priced.

\paragraph{Pricing the probe.}
A probe's worth is entirely informational: it yields an observation
$Z \in \{\text{bite}, \text{no-bite}\}$ that sharpens the belief and therefore the
next decision. We price it by the expected value of sample information
(EVSI)~\cite{howard1966information}:
%
\begin{equation}
V(p) \;=\; \min_a \mathbb{E}[C(a) \mid p] \;-\; \mathbb{E}_Z\!\left[\min_a
\mathbb{E}[C(a) \mid p'(Z)]\right],
\end{equation}
%
where $\beta_a = P(\text{bite} \mid \text{hostile})$ and $\beta_b = P(\text{bite}
\mid \text{benign})$ are the calibrated rates of Section~\ref{sec:calibration},
$P(\text{bite} \mid p) = p\beta_a + (1-p)\beta_b$, and the posteriors follow by
Bayes' rule. Both inner minima are over the same three affine lines, so $V$
evaluates in closed form at any belief and nothing requires simulation. The rule
prices bait at its immediate cost \emph{net of} the information it buys:
$\text{eff}(\text{bait}) = 1 + 24p - V(p)$.

We claim no new mathematics. EVSI is textbook and the properties below are
elementary. The contribution is the application: a response-side probe \emph{is} a
sample-information purchase, and pricing it as one turns the middle option from a
heuristic into a derived action.

\paragraph{The band is bounded on both sides by construction.}
$\min_a \mathbb{E}[C(a) \mid p]$ is a minimum of affine functions and hence concave,
so Jensen's inequality gives $V(p) \geq 0$: information never hurts. And $V(0) =
V(1) = 0$, since at certainty no observation can change the decision. The band
therefore cannot swallow the probability range or be widened by tuning; its edges
are pinned by the cost geometry. Both are enforced as runtime invariants rather than
argued informally.

Bait is chosen where $V(p) > 1 - p$ and the belief is still below the divert line.
With the frozen table and the calibrated library, taking the most informative probe
applicable to the response in hand, the derived bands are
%
\begin{equation}
\text{PASS } p < 0.0647, \qquad
\text{BAIT } 0.0647 \leq p < 0.8793, \qquad
\text{DIVERT } p \geq 0.8793 .
\end{equation}
%
The two edges are set by two different probes, which is what taking a maximum over
the library means. Note the direction of the upper edge: offering a probe
\emph{raises} the divert threshold from $0.816$ to $0.879$, so a system that can
probe is more reluctant to divert an honest user than the same system without one.
That falls out of the arithmetic; it was not a design goal.

\paragraph{One term the textbook derivation does not imply.}
EVSI prices a probe as though its outcome were fresh information, true the first
time and false the tenth. Valuing every exposure at full price lets a bait-aware
attacker who never bites be baited indefinitely rather than diverted, caught by the
passive baseline and not by the full system. Decaying the value geometrically in the
number of unrewarded showings, $(1-\beta_a)^k$, restores the guarantee: as the probe
stops answering, the rule falls back to the two-action policy.

\paragraph{The conclusions survive both estimated inputs.}
Sweeping $\beta_a$ across $[0.05, 0.99]$, and the cost table's one judgement across
$0.5\times$ to $128\times$ its frozen value, the band is non-empty at every point
and the divert threshold never falls below the cost-only boundary. What moves is
\emph{where} the boundaries sit, never \emph{whether} the third action exists, nor
whether a probe could make the system divert earlier than cost accounting alone (it
cannot). The frozen table is a choice about conservatism, not one that manufactures
the result.

\begin{table}[t]
\centering
\caption{The frozen cost matrix. Hashed before any data was collected, so no
threshold in the derivation can have been tuned to the results.}
\label{tab:costs}
\begin{tabular}{lrrr}
\toprule
 & pass & bait & divert \\
\midrule
client benign  & 0  & 1  & 200 \\
client hostile & 25 & 25 & $-20$ \\
\bottomrule
\end{tabular}
\end{table}

% decision-bands.pdf removed for space: Eq. (3) states the derived edges exactly and
% the paragraph above gives the cost-only boundary, so the figure illustrated a
% result already stated numerically. The paper is text-and-table only.

\subsection{Calibrating the Probe}
\label{sec:calibration}

The rule needs one empirical input: how much a bite tells you. Set it by hand to make
the system work and the derived band is a fiction dressed as a derivation. The
evidence a bite carries is a likelihood ratio $\beta_a/\beta_b$, applied in log-odds
so the update is a single addition. A bite moves the malice score sharply up; a
declined probe nudges it gently down. That second ratio matters as much as the first,
because a system that only revises suspicion upward will eventually divert someone
for browsing slowly. For the probe carrying most of our results, 138 of 183 hostile
sessions shown it took it and \textbf{0 of 74 benign ones did}, giving $\beta_a =
0.753$ and $\beta_b = 0.0067$: a bite worth $+4.7$ in log-odds, a refusal $-1.4$.

These rates have no legitimate source in the obvious data. The training corpus
predates the bait library and contains no bites, and the held-out corpus is the test
set~\cite{arp2022dos}. We therefore add a dedicated \emph{calibration} round, run
after the baits exist and frozen before evaluation, whose only permitted use is
estimating probe effectiveness; the round vocabulary is enforced in the record schema
so one round cannot quietly become another. The benign rate is measured rather than
assumed. Setting $\beta_b = 0$ is tempting, since bait is invisible to a real
browser, but it makes the ratio infinite and assumes away the safety property we
exist to measure. Every benign count is in fact zero, so rates are posterior means
under a Jeffreys prior, capping the ratio near 112 rather than claiming a certainty
seventy-four sessions cannot support.

Four of the five deployed baits are measured ($\beta_a$ from 0.563 to 0.794); one
keeps a prior, is flagged as such, and is almost never selected here. A sixth was
\emph{withdrawn}: the target authenticates over HTML, so it reached zero sessions
while its $\beta_a$ remained a number we invented before any data existed, and an
unmeasured value in a measured library is what this section rules out. One caution
worth recording: a calibration harness can be wrong in a way that looks like a
result. Ours understated one probe roughly sevenfold, and nothing raised an error ---
the library stayed internally consistent and the bands looked plausible.

\subsection{Keeping the Decoy Consistent}
\label{sec:consistency}

A decoy that convinces on the first request and contradicts itself on the third has
announced that it is a trap. The hard part is consistency over an engagement:
answering the same query identically twice, keeping a value stable when reached by a
different route, returning what was written, and resolving references between
objects. A fact notebook meets all four, sitting between the decoy's generators and
its responses, recording each value the first time it is needed and serving that
value thereafter. The generator can be anything, since the notebook only requires
that a value never change once produced.

Measuring this needs an adversary, because a decoy asked one question twice passes
almost by accident. A fuzzer interrogates it for the same fact through different
routes, orders and repetitions. With the notebook the contradiction rate over
\textbf{286 adversarial probes is zero}; without it, \textbf{100\%}. The same split
holds against a local language model, so the guarantee belongs to the notebook
rather than to whatever writes the content.

\paragraph{Self-consistency is the wrong estimand.}
Every probe above asks the decoy about the decoy. That is right for the notebook but
not what an attacker tests. A diverted session remembers the \emph{real}
application: it read object ids before being moved and can read them again after, so
a decoy can score zero against itself while contradicting everything the attacker
already saw.

Adversarial re-testing found four tells there, each surfaced by attacking the fix for
the one before: a re-read of the same path returned different values; an aggregate
endpoint named a person differently from their profile page; an id that returned
\emph{not found} before the divert existed afterwards; and a dashboard claimed
ownership of a record the attacker knew was someone else's. All are reachable by an
attacker who simply remembers, and none is visible to a within-decoy metric. So we
report a second measurement: a harness reads records and profiles on the target,
triggers a divert, re-reads the same ids, and counts every field that disagrees
across the boundary. The proxy carries a session's pre-divert observations forward,
including the \emph{absences} it confirmed, so already-seen state survives the move
while anything never looked at stays fabricated. Over 25 diverted sessions and 4,600
compared fields the rate is \textbf{zero}, and the measurement is not vacuous:
disabling the carry-forward reproduces the tells at \textbf{90.8\%}. We stop short of
claiming the class closed, since each tell was found only by attacking the last fix.

% ================================================================= 4 results 1100w
\section{Results}
\label{sec:results}

Section~\ref{sec:rule} makes a prediction that can fail: a priced probe should help
where the passive classifier is genuinely uncertain, and nowhere else. A probe
raising recall everywhere would indicate a weak meter rather than correct pricing,
and one raising recall by diverting honest users would be worthless whatever the
arithmetic said. So we ask, in order: does the probe \emph{cause} more attackers to
be caught, where does the gain live, and what does it cost the benign side?

Three arms differ only in the code path selected by one mode flag against a single
frozen model: \textbf{B1}, a signature WAF; \textbf{B2}, the passive dual meter;
\textbf{B4}, the full system. Each ran over \textbf{99 independent seeded traffic
draws} of 120 attack and 80 benign sessions, giving 11,880 attack and 7,920 benign
sessions per arm. Within a seed every arm sees byte-identical traffic, so each attack
session forms a matched pair. Throughout, \emph{diverted} means the policy decided to
divert, counted identically in every arm, so recall is comparable by construction.
The primary comparison was fixed before the runs.

\subsection{The Causal Estimate}

B2 against B4 compares two whole systems differing in more than the probe. What
isolates it is a \textbf{randomised holdout inside the treated arm}: whenever the
policy decides to bait, a pseudo-random draw over the session id withholds the probe
from about one session in ten. Withheld sessions sit at the same belief, under the
same policy, in the same band, and are simply passed instead of probed.

\begin{table}[t]
\centering
\caption{The randomised holdout. Both groups reached the probing band under the
same policy; only the probe differs.}
\label{tab:holdout}
\begin{tabular}{lrrr}
\toprule
group & $n$ & diverted & divert rate \\
\midrule
baited (policy)     & 10,643 & 10,112 & \textbf{0.950} \\
withheld (holdout)  &  1,237 &  1,089 & \textbf{0.880} \\
\bottomrule
\end{tabular}
\end{table}

The effect is \textbf{+0.070}, bootstrap 95\% CI \textbf{[+0.052, +0.088]}, odds
ratio 2.59, Fisher exact \textbf{$p = 3.4 \times 10^{-19}$}. Probing causes
diversions that would not otherwise have happened, at the same belief state. A
randomised design is only worth the name if the draw balanced, so we check: the
largest share difference across the five subcategories is 1.2 percentage points, and
a chi-square of the withheld mix against the treated gives 2.58 on four degrees of
freedom, well inside 9.49.

\subsection{Baselines}

\begin{table}[t]
\centering
\caption{Attack recall and benign cost over 99 paired seeds. B2 and B4 intervals do
not overlap.}
\label{tab:baselines}
\setlength{\tabcolsep}{10pt}
\begin{tabular}{@{}l l c c@{}}
\toprule
arm & recall (Wilson 95\% CI) & per-seed sd & benign diverted \\
\midrule
B1 signature WAF & 0.366 [0.358, 0.375] & 0.020 & 0 / 7,920 \\
B2 passive       & 0.889 [0.883, 0.894] & 0.024 & 4 / 7,920 \\
\textbf{B4 full} & \textbf{0.943 [0.939, 0.947]} & 0.018 & \textbf{0 / 7,920} \\
\bottomrule
\end{tabular}
\end{table}

Of 11,880 matched attack sessions, 10,841 are concordant, \textbf{842 are caught by
B4 alone and 197 by B2 alone}, giving McNemar $p = 1.9 \times 10^{-95}$; B4 leads in
\textbf{99 of 99 seeds}, which rules out a lucky draw. The 197 going the other way
are not noise. They are \emph{deferrals}: sessions whose belief landed between the
cost-only boundary at $0.816$ and the derived divert edge at $0.879$, where B2
diverts at once and B4 buys information first. Some end without the probe ever being
answered. That is the visible price of the third action.

The two zeros in the benign column invite suspicion, so it is worth saying what
produced them. B1's is unremarkable: its patterns need SQL syntax context that honest
traffic does not supply. B4's is the interesting one, because it is not merely as
safe as the passive arm but safer, by the mechanism Section~\ref{sec:rule} predicts.
All four benign sessions B2 wrongly diverted landed in the probing band under B4.
Each was probed rather than diverted, none bit, and the \emph{absence} of a bite held
the belief below $0.879$ where B2's boundary of $0.816$ had already committed.
Raising the divert threshold is not a safety concession bolted on afterwards; it is
what the pricing implies, and here it recovered every false positive the passive
detector produced.

B1 is a fair reference rather than a straw man, but a reader is entitled to distrust
a baseline its authors wrote. So we replayed the same traffic through \textbf{OWASP
ModSecurity CRS} at every paranoia level it defines.

\begin{table}[t]
\centering
\caption{OWASP ModSecurity CRS on identical traffic. Perfect recall is reachable
only at a setting no operator can run.}
\label{tab:crs}
\begin{tabular}{lll}
\toprule
CRS paranoia & attack recall (95\% CI) & benign sessions blocked \\
\midrule
1 (default) & 0.353 [0.305, 0.404] & \textbf{0 / 240} \\
2           & 0.544 [0.493, 0.595] & \textbf{0 / 240} \\
3           & 0.544 [0.493, 0.595] & \textbf{0 / 240} \\
4 (maximum) & \textbf{1.000} [0.989, 1.000] & \textbf{72 / 240 (30\%)} \\
\bottomrule
\end{tabular}
\end{table}

At its default setting the real ruleset scores 0.353 against B1's 0.366 with heavily
overlapping intervals, so B1 is not a weakened stand-in. The ceiling is not a tuning
question either: CRS reaches perfect recall only by blocking \textbf{thirty per cent
of legitimate sessions}, and across every setting that leaves benign traffic alone it
tops out near 0.54, against 0.943 at zero benign diversions here. The comparison is
generous to CRS in one respect: it is judged on blocking any request of a session,
while our arms are judged on diverting the session, and blocking is easier.

\subsection{Where the Probe Acts}

Splitting the matched pairs by subcategory is the sharpest test of the theory, which
says the gain must be \emph{concentrated}. It is. The entire net gain sits in
UI-based scattered object access: recall $0.540 \rightarrow 0.823$, a gain of
\textbf{+0.282} over 1,980 pairs at a bite rate of 0.620, with a paired McNemar on
that subcategory alone of $b = 639$ against $c = 80$ ($p < 10^{-4}$). Elsewhere the
probe is neutral: $+0.022$ on obfuscated SQLi, and zero on stealth SQLi, API-side
object access and credential spray. That is the one category a signature cannot see
and the one where the passive meter is undecided, exactly where
Section~\ref{sec:rule} says information is worth buying. Stealth SQLi bites half the
time and does not move at all: the information was bought and proved unnecessary.

One category cannot reach the probe at all, and it is half of what B4 still misses.
Obfuscated SQLi bites exactly zero times despite nine in ten of its sessions being
shown a probe, and accounts for 405 of B4's 799 misses. The cause lies in the
attacker model, not the defence: those profiles fire payloads and never read what
comes back, so a response-side probe cannot reach them. Ablations, an
adaptive-adversary sweep, a transfer check against a second application, and runs
against sqlmap, ghauri, OWASP ZAP and Wapiti are reported in the artefact.

% holdout-effect.pdf removed: Table~\ref{tab:holdout} already carries both rates and
% the n, so the figure restated it and cost most of a page at 12-page limit.

% ============================================================== 5 discussion 450w
\section{Discussion}
\label{sec:discussion}

The results support the mechanism rather than a headline number. A probe lifting
recall everywhere would suggest a weak meter; instead the gain of $+0.282$ is
confined to the one subcategory where the passive detector is undecided and is
roughly zero in the four where it is not. That is the pricing argument doing what it
predicts: information is worth buying exactly where the decision is in doubt. The
holdout separates that effect from every other difference between two systems, and
the raised divert edge explains why the safest arm is also the most accurate.

\paragraph{Where the approach does not win.}
Deriving the band edges did not beat setting them by hand on expected cost, and we
report the measured reason rather than the claim. The belief those edges apply to
takes essentially two distinct values on this system, so four decimal places of
derived precision are not doing the work they appear to, and every configuration
beating the derived pair on cost does so by diverting benign users. Among
configurations that divert none, the derived edges are the best available. That is
narrower than we set out to claim and more useful: it says what the derivation buys,
a placement clearing the benign belief distribution by construction, rather than
asserting a superiority the data does not support. A related gap is structural. The
rule is derived for a single decision but deployed as a first-crossing test over a
session, and here that mismatch happens to cancel against the meter's
miscalibration. A rule that is right for compensating reasons is a different object
from a rule that is right.

\paragraph{Limitations.}
Four, in order of severity. \emph{All traffic is synthetic}, so every rate describes
this distribution and would need re-measuring on a real one; the derivation and the
holdout design do not depend on the traffic being real, but the magnitudes do.
\emph{A single tuned application}: placing the frozen model in front of a
structurally different second one shows the app-agnostic features fire correctly,
but the app-specific probes would need re-pointing, and a full evaluation with a
matched decoy there remains undone. \emph{The deception is not validated by human
participants}; whether a human attacker \emph{feels} something is off is unmeasured,
and it is the single most valuable thing left to measure. \emph{$\beta_a$ is a
property of an attacker model we chose.} We narrow this by driving autonomous
language-model agents told nothing about the probes, whose bite rate rises with
capability and whose largest lands inside the range we assumed. An agent is still
not a person.

% ============================================================== 6 conclusion 120w
\section{Conclusion}

Passive web defences force a choice between acting early on weak evidence and acting
late on strong evidence. That choice is avoidable: a defender can \emph{manufacture}
evidence by planting an invisible, inert probe, and when to do so follows from the
expected value of the information it buys. Priced this way, probing is the
cost-optimal action over a band that does not exist under cost accounting alone, so
it cannot be a tuned threshold. The evaluation supports the mechanism rather than a
number: the gain appears only where the classifier is genuinely unsure, a randomised
holdout attributes it causally to the probe, and no legitimate session was diverted.
What remains is real traffic, and a real attacker.

% Squeeze the inter-entry leading in the reference list. Purely typographic: no
% entry is dropped or shortened, the list simply stops wasting a page on padding.
\let\oldthebibliography\thebibliography
\renewcommand{\thebibliography}[1]{%
  \oldthebibliography{#1}%
  \setlength{\itemsep}{0pt}%
  \setlength{\parskip}{0pt}%
}
\bibliographystyle{splncs04}
\bibliography{refs}

\end{document}
